\documentclass{article}

\usepackage[utf8]{inputenc}
\usepackage[T1]{fontenc}
\usepackage{helvet}
\usepackage{geometry}
\usepackage{xcolor}
\usepackage{eso-pic}
\usepackage{fancyhdr}
\usepackage{parskip}
\usepackage{microtype}
\usepackage[english, ukrainian]{babel}
\usepackage{url}
\usepackage{amsmath}
\usepackage{amssymb}
\usepackage{amsthm}
\usepackage{longtable}
\usepackage{booktabs}
\usepackage{hyperref}
\usepackage{titlesec}

\definecolor{nato}{HTML}{293757}
\definecolor{footergray}{gray}{0.65}

\geometry{a4paper,left=3cm,right=3cm,top=3cm,bottom=3cm}
\renewcommand{\familydefault}{\sfdefault}
\setcounter{secnumdepth}{3}

\titleformat{\section}
  {\color{nato}\Huge\bfseries}{\thesection.}{1em}{}
\titlespacing*{\section}{0pt}{30pt}{12pt}

\titleformat{\subsection}
  {\color{nato}\LARGE\bfseries}{\thesubsection.}{1em}{}
\titlespacing*{\subsection}{0pt}{20pt}{8pt}

\titleformat{\subsubsection}
  {\color{nato}\Large\bfseries}{\thesubsubsection.}{1em}{}
\titlespacing*{\subsubsection}{0pt}{10pt}{6pt}

\fancyhf{}
\fancyhead[R]{\small\color{footergray} 27 червня 2026}
\fancyfoot[L]{\small\color{footergray} Техноробочий проект ERP/1. Модуль Документи}
\fancyfoot[R]{\small\color{footergray}\thepage}

\newcommand\HUGE[1]{\fontsize{40}{40}\selectfont #1}

\renewcommand{\headrulewidth}{0pt}
\renewcommand{\footrulewidth}{0.4pt}
\renewcommand{\footrule}{\color{footergray}\hrule width \textwidth height 0.4pt}

\begin{document}

\pagestyle{fancy}

\begin{titlepage}
  \AddToShipoutPictureBG*{\AtPageLowerLeft{\color{nato}\rule{\paperwidth}{\paperheight}}}
  \color{white}\thispagestyle{empty}\vspace*{3cm}\centering
  {\Huge  \textbf{ERP/1: Документи} \par} \vspace{1.5cm}
  {\HUGE  \textbf{Техноробочий проект} \par} \vspace{1.5cm}
  {\LARGE \textbf{на розробку та впровадження підсистеми електронного документообігу та обміну СЕВ ОВВ} \par} \vspace{1.5cm}
  {\large Відповідно до ДСТУ 34.201-89 та Постанови КМУ № 205 від 21.02.2025 \par}
  {\large Формалізація вимог за стандартом ISO/IEC/IEEE 29148:2018 \par}
  \vfill
  {\large 2026 \copyright\ Системи Електронного Урядування \par}
\end{titlepage}

\newpage

\begin{center}
  {\Huge\color{nato}\bfseries Зміст\par}
\end{center}
\vspace{40pt}
\tableofcontents

\newpage

\section{Загальні відомості про засіб інформатизації}

\subsection{Найменування та підстави для розробки}
\textbf{Найменування засобу інформатизації:} Підсистема «ERP/1: Документи» (далі --- Модуль «Документи» або СЕД).
\\
\textbf{Відомості про замовника:} ТОВ "Системи Електронного Урядування".
\\
\textbf{Підстави для виконання робіт:}
\begin{itemize}
    \item Закон України «Про електронний документообіг»;
    \item Постанова Кабінету Міністрів України від 17 січня 2018 року № 55 «Деякі питання документування управлінської діяльності»;
    \item Постанова Кабінету Міністрів України від 21 лютого 2025 року № 205 «Деякі питання створення, адміністрування та забезпечення функціонування засобу інформатизації».
\end{itemize}

\subsection{Призначення та цілі створення засобу}
Модуль «Документи» призначений для організації проходження та автоматизації життєвого циклу електронних документів, візування, підписання КЕП, контролю за виконанням завдань, а також міжвідомчого поштового обміну за стандартами X.420 / X.422.3.
\\
\textbf{Цілі створення засобу:}
\begin{itemize}
    \item Автоматизація вхідної, вихідної та внутрішньої кореспонденції;
    \item Контроль виконавської дисципліни за допомогою резолюцій керівництва;
    \item Підвищення прозорості та швидкості погодження проектів документів;
    \item Інтеграція з СЕВ ОВВ через систему електронної взаємодії;
    \item Надійне та масштабоване збереження документів у СУБД Mnesia та RocksDB.
\end{itemize}

\subsection{План-графік виконання етапів робіт}
\begin{enumerate}
    \item \textbf{Етап 1: Розробка архітектурних рішень та інтерфейсів} --- Створення карти продукту в N2O.DEV, дизайн-прототипів у Figma. (1.5 місяці).
    \item \textbf{Етап 2: Реалізація FSM бізнес-процесів} --- Написання Elixir BPE специфікацій та ABAC правил доступу. (2 місяці).
    \item \textbf{Етап 3: Створення сховища даних та КЕП сервісів} --- Опис рекорді Mnesia, RocksDB схем, інтеграція з бібліотеками валідації КЕП/КЕПч. (1.5 місяці).
    \item \textbf{Етап 4: Інтеграція з СЕВ ОВВ} --- Налаштування шлюзу обміну поштовими пакетами X.420. (2 місяці).
    \item \textbf{Етап 5: Тестування та приймально-здавальні випробування} --- Проведення Unit, UAT та тестів надійності. (1 місяць).
\end{enumerate}

\newpage
\section{Відомості про робочий процес та умови експлуатації}

\subsection{Опис бізнес-процесів (BPMN / FSM)}

\subsubsection{Процес обробки вхідних документів (INPUT)}
\textbf{Опис кроків та станів (BPE FSM Transitions):}
\begin{enumerate}
    \item \textbf{Created (Чернетка):} Реєстратор створює картку вхідного документа. Перехід на наступний крок можливий лише після накладання системного КЕП/КЕПч та заповнення обов'язкових полів.
    \item \textbf{Registration (Реєстрація):} Присвоєння вхідного реєстраційного номера за формулою журналу.
    \item \textbf{gwNeedsDetermination (Умовний шлюз):} Автоматична перевірка: якщо документ терміновий (\texttt{urgent == true}), він перенаправляється на етап \texttt{UrgentConsideration}, інакше --- на \texttt{InitialConsideration}.
    \item \textbf{InitialConsideration (Попередній розгляд):} Реєстратор визначає виконавця або керівника для накладання резолюції.
    \item \textbf{Resolution (Резолюція):} Керівник накладає доручення (задачу).
    \item \textbf{Implementation (Виконання):} Створюється бізнес-задача \texttt{bizTask}.
    \item \textbf{Completed (Виконано):} Документ списується в архів.
\end{enumerate}

\subsubsection{Процес обробки вихідних документів (OUTPUT)}
\textbf{Опис кроків та умов переходів:}
\begin{enumerate}
    \item \textbf{Draft (Чернетка):} Створення проекту документа виконавцем, додавання файлів (\texttt{attachments}).
    \item \textbf{Coordination (Візування):} Процес внутрішнього погодження. Якщо хоча б один візант відхиляє проект (\texttt{reject == true}), процес повертається на крок \texttt{Draft} з коментарем.
    \item \textbf{Signing (Підписання):} Передача керівнику для накладання КЕП.
    \item \textbf{Registration (Реєстрація):} Автоматична реєстрація вихідного номера.
    \item \textbf{Dispatch (Відправка):} Передача через СЕВ ОВВ (конверт X.420).
\end{enumerate}

\subsubsection{Процес внутрішніх документів (INTERNAL)}
\begin{enumerate}
    \item \textbf{Draft:} Створення внутрішньої записки чи проекту наказу.
    \item \textbf{Coordination:} Візування посадовими особами.
    \item \textbf{Signing:} Підписання КЕП керівником.
    \item \textbf{Registration:} Реєстрація в журналі внутрішніх документів.
    \item \textbf{Distribution:} Автоматична розсилка адресатам усередині установи для ознайомлення чи виконання.
\end{enumerate}

\subsubsection{Організаційно-розпорядчі документи (ORG)}
\begin{enumerate}
    \item \textbf{Draft:} Створення проекту наказу/розпорядження.
    \item \textbf{JuridicalApproval:} Обов'язковий юридичний аналіз юридичним департаментом.
    \item \textbf{Coordination:} Погодження членами колегії чи заступниками.
    \item \textbf{Signing:} Підписання КЕП головою установи.
    \item \textbf{Registration:} Реєстрація та розсилання виконавцям із автоматичним створенням контрольних задач (\texttt{ControlTask}).
\end{enumerate}

\subsubsection{Контроль виконання доручень (CONTROLTASK)}
\begin{enumerate}
    \item \textbf{ControlStart:} Автоматичне взяття задачі на контроль.
    \item \textbf{Monitoring:} Регулярний аналіз прогресу виконання контролюючим органом.
    \item \textbf{Warning:} Автоматичне сповіщення виконавця за 3 дні до закінчення терміну виконання.
    \item \textbf{Completed:} Зняття з контролю після погодження звіту контролером.
\end{enumerate}

\subsubsection{Виконання задач співвиконавцями (EXECTASK)}
\begin{enumerate}
    \item \textbf{TaskReceived:} Отримання завдання виконавцем.
    \item \textbf{CoordinationWithChief:} Підготовка проміжних рішень співвиконавцями.
    \item \textbf{ReportPreparation:} Написання звіту про виконання та прикріплення відповіді.
    \item \textbf{Submission:} Передача результату головному виконавцю.
\end{enumerate}

\subsubsection{Розгляд звернень громадян (ADMISSION)}
\begin{enumerate}
    \item \textbf{AdmissionReceived:} Реєстрація звернення згідно з КОАТУУ та категоріями заявників.
    \item \textbf{Consideration:} Попередній розгляд та резолюція (термін розгляду до 30 календарних днів).
    \item \textbf{ResponseDrafting:} Підготовка відповіді заявнику.
    \item \textbf{Dispatch:} Відправка відповіді КЕП та списання справи.
\end{enumerate}

\subsubsection{Додаткові бізнес-процеси (ADDITIONAL)}
\begin{enumerate}
    \item \textbf{Start:} Реєстрація додаткових запитів чи адвокатських запитів (термін виконання --- 5 робочих днів).
    \item \textbf{Execution:} Асинхронне отримання довідок від підрозділів.
    \item \textbf{Response:} Відправка результату.
\end{enumerate}

\subsection{Опис ролей та прав доступу (ABAC)}
Автентифікація та авторизація виконуються на базі ABAC атрибутів:
\begin{itemize}
    \item \textbf{Правило 1 (Реєстратор):}
    $$\text{subject.role contains :registrar} \land \text{object.entity\_type == :inputOrder}$$
    $$\land\ \text{object.status == :draft} \rightarrow \text{Permit: create, read, write, register.}$$
    \item \textbf{Правило 2 (Виконавець):}
    $$\text{subject.role contains :executor} \land \text{object.entity\_type == :bizTask}$$
    $$\land\ (\text{object.executor == subject.id} \lor \text{object.subexecutors contains subject.id})$$
    $$\rightarrow \text{Permit: read, write (виконання, прикріплення відповідей).}$$
    \item \textbf{Правило 3 (Керівник):}
    $$\text{subject.role contains :head} \land \text{object.entity\_type == :inputOrder}$$
    $$\rightarrow \text{Permit: read, write (накладання резолюцій, підписання КЕП).}$$
    \item \textbf{Правило 4 (Помічник/Асистент):}
    $$\text{subject.role contains :assistant} \land \text{subject.substitute == object.head}$$
    $$\rightarrow \text{Permit: read, write (підготовка проектів резолюцій/підписів).}$$
\end{itemize}

\subsection{Умови експлуатації та технологічний стек}
\begin{itemize}
    \item Програмна платформа: Erlang/OTP версії 26 та Elixir версії 1.15.
    \item База даних: СУБД Mnesia для оперативних таблиць, RocksDB як сховище KVS для довгострокового зберігання історичних ланцюжків та бінарних об'єктів.
    \item Інтерфейс користувача: Web-клієнт на базі протоколів N2O та Nitro UI.
\end{itemize}

\newpage
\section{Інформаційне та програмне забезпечення}

\subsection{Специфікація ERP реквізитної інформації (Records)}

Нижче наведено детальні структури даних з include-файлів бібліотеки \texttt{schema} для збереження документів та довідників СЕД.

\subsubsection{Вхідний документ (inputOrder)}
\begin{verbatim}
-record(inputOrder, { id= [] :: binary(),
                      guid= [] :: list(),
                      urgent=[] :: [] | boolean(),
                      date=[] :: [] | calendar:datetime(),
                      xml = [] :: [] | binary(),
                      hash = [] :: binary() | list(),
                      signature = [] :: binary() | list(),
                      document_type= <<"Вхідний документ"/utf8>> :: term(),
                      nomenclature=[] :: term(),
                      category=[] :: [] | list(),
                      corr=[] :: [] | list(),
                      branch=[] :: [] | list(),
                      head=[] :: [] | #'Person'{},
                      signed=[] :: [] | list(),
                      number_out= [] :: list(),
                      date_out=[] :: calendar:datetime(),
                      content= [] :: list(),
                      note=[] :: list(),
                      dueDate=[] :: [] | calendar:datetime(),
                      dueDateTime = [] :: list(),
                      to=[] :: [] | #'Person'{},
                      main_sheets=[] :: [] | list(),
                      add_sheets=[] :: [] | list(),
                      device = [] :: [] | list(),
                      registered_by=[] :: [] | #'Person'{},
                      coordination=[] :: list(),
                      topic=[] :: [] | binary(),
                      attachments = [] :: list(term()),
                      seq_id = [] :: list(),
                      reject_comment = [] :: binary(),
                      created_by = [] :: term(),
                      created = [] :: term(),
                      modified_by = [] :: term(),
                      modified = [] :: term(),
                      bizTask_initiator = [] :: term(),
                      toAttention = #assistantMark{},
                      project = #project{} :: #project{}
 }).
\end{verbatim}

\subsubsection{Вихідний документ (outputOrder)}
\begin{verbatim}
-record(outputOrder, { id= [] :: binary(),
                       guid= [] :: list(),
                       urgent=[] :: [] | boolean(),
                       date=[] :: [] | calendar:datetime(),
                       xml = [] :: [] | binary(),
                       hash = [] :: binary() | list(),
                       signature = [] :: binary() | list(),
                       document_type = [] :: tuple(),
                       nomenclature=[] :: term(),
                       purpose=[] :: [] | list(),
                       purpose_sev=[] :: [] | list(),
                       corr=[] :: [] | list(),
                       deputy=[] :: [] | list(),
                       sign=[] :: [] | list(),
                       executor = [] :: [] | #'Person'{},
                       branch = [] :: list(),
                       head = [] :: [] | #'Person'{},
                       content= [] :: list(),
                       approvers = [] :: list(#'Person'{}),
                       signatory = [] :: list(#'Person'{}),
                       send_type = [] :: list(),
                       sent = #sendInfo{} :: #sendInfo{},
                       registered_by=[] :: [] | #'Person'{},
                       date_send=[] :: [] | calendar:datetime(),
                       note=[] :: list(),
                       attachments = [] :: list(term()),
                       project = #project{} :: #project{},
                       seq_id = [] :: list(),
                       created_by = [] :: term(),
                       created = [] :: term(),
                       modified_by = [] :: term(),
                       modified = [] :: term(),
                       template = [] :: term(),
                       assistantProcessed = #assistantMark{},
                       toAttention = #assistantMark{}
              }).
\end{verbatim}

\subsubsection{Внутрішній документ (internalDoc)}
\begin{verbatim}
-record(internalDoc, { id= [] :: binary(),
                       guid= [] :: list(),
                       urgent=[] :: [] | boolean(),
                       date=[] :: [] | calendar:datetime(),
                       xml = [] :: [] | binary(),
                       hash = [] :: binary() | list(),
                       signature = [] :: binary() | list(),
                       document_type = [] :: tuple(),
                       nomenclature=[] :: term(),
                       executor=[] :: [] | #'Person'{},
                       approvers = [] :: list(#'Person'{}),
                       branch=[] :: [] | list(),
                       head=[] :: [] | #'Person'{},
                       content=[] :: term(),
                       signatory=[] :: [] | list(#'Person'{}),
                       target=[] :: [] | #'Person'{},
                       note=[] :: list(),
                       attachments = [] :: list(term()),
                       registered_by=[] :: [] | #'Person'{},
                       project = #project{} :: #project{},
                       seq_id = [] :: list(),
                       created_by = [] :: term(),
                       created = [] :: term(),
                       modified_by = [] :: term(),
                       modified = [] :: term(),
                       template = [] :: term(),
                       bizTask_initiator = [] :: term(),
                       toAttention = #assistantMark{}
                      }).
\end{verbatim}

\subsubsection{Бізнес-задача / Резолюція (bizTask)}
\begin{verbatim}
-record(bizTask, { id= [] :: binary(),
                   task_type = [] :: [] | list(),
                   initiator = [] :: [] | #'Person'{},
                   point=[] :: [] | list(),
                   content = [] :: list(),
                   document_type= <<"Задача"/utf8>> :: term(),
                   nomenclature=[] :: term(),
                   sign = [] :: [] | boolean(),
                   executor=[] :: [] | #'Person'{},
                   subexecutors = [] :: [] | list(#'Person'{}),
                   notify = [] :: [] | list(#'Person'{}),
                   fromDate = [] :: [] | calendar:datetime(),
                   dueDate = [] :: [] | calendar:datetime(),
                   dueDate_subexecutors = [] :: [] | calendar:datetime(),
                   priority = [] :: term(),
                   sent = #sendInfo{},
                   taskTime = [] :: integer(),
                   control=[] :: [] | boolean(),
                   period=[] :: [] | list(),
                   date_start=[] :: [] | calendar:datetime(),
                   date_finish=[] :: [] | calendar:datetime(),
                   date_control=[] :: [] | calendar:datetime(),
                   control_by=[] :: [] | list(#'Person'{}),
                   comment=[] :: [] | list(),
                   remove= false :: boolean(),
                   registered_by=[] :: [] | #'Person'{},
                   removed_by=[] :: [] | #'Person'{},
                   removed_date=[] :: [] | calendar:datetime(),
                   date_complete=[] :: [] | calendar:datetime(),
                   controller_comment=[] :: [] | list(),
                   complete_comment=[] :: [] | list(),
                   progress = [] :: [] | list(),
                   attachments = [] :: list(term()),
                   seq_id = [] :: list(),
                   parent = [] :: term(),
                   project = [] :: term(),
                   created_by = [] :: term(),
                   created = [] :: term(),
                   modified_by = [] :: term(),
                   modified = [] :: term(),
                   notifyTaskExist = false :: false | boolean(),
                   toAttention = #assistantMark{}
                 }).
\end{verbatim}

\subsubsection{Співробітник (Employee)}
\begin{verbatim}
-record('Employee',    { id          = kvs:seq([],[]) :: term(),
                         person      = [] :: [] | #'Person'{},
                         org         = [] :: [] | #'Organization'{},
                         branch      = [] :: #'Branch'{} | list(),
                         code        = [],
                         surname     = [],
                         name        = [],
                         middle_name = [],
                         sex         = [],
                         birthday    = [],
                         hired       = [],
                         fired       = [],
                         inn         = [],
                         position    = [] :: list(),
                         number      = [],
                         group       = [] :: list(),
                         role        = [] :: list(atom()),
                         substitute  = [] :: [] | #'Person'{},
                         substitutes = [] :: list(),
                         rightsDelegation = [] :: list() | term(),
                         phone       = [] :: list() | binary(),
                         mail        = [] :: list(),
                         status      = [] :: term(),
                         type        = [] :: term() }).
\end{verbatim}

\subsubsection{Організація (Organization)}
\begin{verbatim}
-record('Organization', { id = [] :: list(),
                          name        = [] :: [] | binary(),
                          details     = [] :: [] | binary(),
                          hq          = "00000" :: term(),
                          orgname     = [] :: [] | binary(),
                          address     = [] :: [] | binary(),
                          house       = [] :: [] | binary(),
                          director    = [] :: [] | term(),
                          position    = [] :: [] | binary(),
                          phones      = [] :: list() | binary(),
                          fax         = [] :: list() | binary(),
                          url         = [] :: [] | string(),
                          code        = [] :: list() | binary(),
                          login       = [] :: [] | binary(),
                          password    = [] :: [] | binary(),
                          location    = [] :: [] | #'Loc'{},
                          type        = [] :: term(),
                          keys        = 50 :: integer(),
                          email       = [] :: [] | binary(),
                          feed        = [] :: [] | binary(),
                          sendType    = [] :: [] | term() }).
\end{verbatim}

\subsubsection{Фізична особа (Person)}
\begin{verbatim}
-record('Person',      { id          = kvs:seq([],[]) :: [] | term(),
                         cn          = [] :: [] | binary(),
                         name        = [] :: [] | binary(),
                         displayName = [] :: [] | binary(),
                         location    = [] :: #'Loc'{},
                         hours       = 00 :: integer(),
                         type        = [] :: term() }).
\end{verbatim}

\subsubsection{Департамент / Філія (Branch)}
\begin{verbatim}
-record('Branch',   { id     = [] :: [] | binary(),
                      name   = [] :: list(),
                      org    = [] :: [] | binary(),
                      parent = [] :: [] | binary(),
                      short  = [] :: list(),
                      head   = [] :: [] | #'Person'{},
                      deputies   = [] :: list(#'Person'{}),
                      type   = [] :: [] | atom(),
                      lvl    = 1 :: integer(),
                      index  = [] :: [] | binary(),
                      loc    = [] :: [] | #'Loc'{} }).
\end{verbatim}

\subsection{Програмне забезпечення (Реалізація)}

\subsubsection{newInputProc.ex (Процес обробки вхідних документів)}

\begingroup
\footnotesize
\begin{verbatim}
defmodule NewInput.Proc do
  import General.Proc
  import Doclink
  require BPE
  require ERP

  def def(), do: BPE.xml("/priv/bpmn/newInputProc.bpmn")
  def auth(_roles), do: true

  def action(BPE.broadcastEvent(), BPE.process() = proc) do
    BPE.result(type: :stop, state: proc)
  end

  def action(BPE.messageEvent(), BPE.process(userStarted: :system) = proc) do
    BPE.result(type: :stop, state: proc)
  end

  def action(BPE.messageEvent(payload: {:next, _}) = r, BPE.process() = proc) do
    NewOutput.Proc.action(r, proc)
  end

  def action(BPE.messageEvent(name: 'DocumentStatistic',
    payload: {executed, executors}), BPE.process(id: pid) = proc) do
    case :bpe.head(pid) do
      BPE.hist(task: BPE.sequenceFlow(source: s, target: t)) ->
        NewOutput.Proc.personal_statistic({:callback, s, t},
          BPE.result(executed: executed,
          state: BPE.process(proc, executors: executors)), :document)
      _ -> []
    end
    BPE.result(type: :stop, state: proc)
  end

  def action(BPE.messageEvent(name: 'RemoveDocumentStatistic',
    payload: {stage, executors}), BPE.process() = proc) do
    NewOutput.Proc.remove_personal_statistic(:executed, :document,
      proc, stage, executors)
    BPE.result(type: :stop, state: proc)
  end

  def action(BPE.messageEvent(name: 'Implementation', payload: initiator),
    BPE.process(id: pid, docs: [doc | _]) = proc) do
    newDoc = :kvs.setfield(doc, :project,
      CRM.KEP.projAction(:execute, :kvs.field(doc, :project), initiator, []))
    BPE.result(state: BPE.process(proc, docs: [newDoc]),
      continue: [BPE.continue(fn: :next, args: [pid])])
  end

  def action(BPE.messageEvent(name: 'Executed',
    sender: BPE.process(id: pid, docs: [doc | _])),
    BPE.process(id: subj) = proc) do
    BPE.result(state: proc, continue: [
      BPE.continue(fn: :generateHistory, type: :spawn, module: :'Elixir.History',
        args: [%HistExt{action: :executed, procId: subj, subPid: pid,
        newValue: doc}]),
      BPE.continue(fn: :next, args: [subj])
    ])
  end

  def action(BPE.messageEvent(name: 'TaskCreated',
    payload: ERP.notifyTask(notify: n)), BPE.process(docs: [doc | _]) = proc) do
    newDoc = :kvs.setfield(doc, :modified_by, n ++
      case :kvs.field(doc, :modified_by) do
        x when is_list(x) -> x; _ -> []
      end)
    BPE.result(type: :stop, state: BPE.process(proc, docs: [newDoc]))
  end

  def action(BPE.messageEvent(), BPE.process() = proc) do
    BPE.result(type: :stop, state: proc)
  end

  def action(BPE.asyncEvent(name: n) = ev, BPE.process() = proc)
    when n in ['Reject', 'gwReject', 'Edit'] do
    NewOutput.Proc.action(ev,  proc)
  end

  def action(BPE.asyncEvent(name: s) = r, state)
    when s in ['RemoveStatus', 'SetStatus'] do
    NewOutput.Proc.action(r, state)
  end

  def action(BPE.asyncEvent(name: 'ImplementationBlock',
    sender: BPE.process(id: pid),
    payload: ERP.bizTask(initiator: i, executor: e,
    subexecutors: s, notify: n, control_by: c)),
    BPE.process(id: subj, docs: [doc | _]) = proc) do
    :bpe.gw_block(pid, 'gwConfirmation', subj)
    newDoc = :kvs.setfield(doc, :modified_by,
      :lists.flatten([e, s, n, c]) ++
      case :kvs.field(doc, :modified_by) do
        x when is_list(x) -> x; _ -> []
      end) |> :kvs.setfield(:bizTask_initiator, i)
    BPE.result(type: :stop, state: BPE.process(proc, docs: [newDoc]))
  end

  def action(BPE.asyncEvent(name: 'ImplementationBlock',
    sender: BPE.process(id: pid)), BPE.process(id: subj) = proc) do
    :bpe.gw_block(pid, 'gwConfirmation', subj)
    BPE.result(type: :stop, state: proc)
  end

  def action(BPE.asyncEvent(name: 'Delete', payload: u),
    BPE.process(id: pid, parent: p, docs: [doc | _], name: n) = proc) do
    u != [] && History.generateHistory(%HistExt{action: :delete_proc,
      procId: p, user: u, docNum: RMK.Control.rowGetName(doc, n)})
    BPE.result(state: proc, continue: [
      BPE.continue(type: :stop),
      BPE.continue(type: :spawn, fn: :delete_process,
        module: :"Elixir.General.Proc", args: [pid])
    ])
  end

  def action(BPE.asyncEvent(name: 'History') = ev, BPE.process() = proc) do
    NewOutput.Proc.action(ev,  proc)
  end

  def action(BPE.asyncEvent(), proc) do
    BPE.result(type: :stop, state: proc)
  end

  def action({:request, _, _} = r,
    BPE.process(userStarted: :system) = proc) do
    BPE.result(type: :stop, state: generalAction(r, proc, proc))
  end

  def action({:request, 'Created', 'Registration'} = r,
    BPE.process(id: pid, docs: [doc | _], parent: p) = proc0) do
    {_, proc} = :bpe.ensure_mon(proc0)
    ERP.project(comment: newComment) = rejectComment(:kvs.field(doc, :project))
    newDoc = :kvs.setfield(doc, :project, ERP.project(comment: newComment))
    :bpe.broadcastEvent(pid, BPE.broadcastEvent(name: 'Delete', sender: proc))
    saveChildProc(pid, p)
    CRM.Search.Common.save_index(doc, [], pid)
    CRM.Search.Common.save_index([:corr, :number_out, :date_out], doc, [], pid)
    newProc = BPE.process(proc, docs: [newDoc])
    :bpe.load(p) != [] && :bpe.messageEvent(p,
      BPE.messageEvent(name: 'Registration', payload: newProc))
    BPE.result(state: generalAction(r, proc, newProc),
      continue: [
        continueGeneral(r, newProc, newDoc,
          if newComment == [] do :created else [] end),
        BPE.continue(type: :stop)
      ])
  end

  def action({:request, 'Registration', 'gwNeedsDetermination'} = r,
    BPE.process(id: pid, executors: e, docs: [doc | _]) = proc) do
    newDoc = :kvs.setfield(doc, :project, rejectComment(:kvs.field(doc, :project)))
    newProc = BPE.process(proc, docs: [newDoc])
    indexUrgent(newProc, newDoc)
    BPE.result(state: generalAction(r, proc, newProc), executed: e,
      continue: [
        continueGeneral(r, newProc, newDoc),
        BPE.continue(fn: :next, args: [pid])
      ])
  end

  def action({:request, 'InitialConsideration', _} = r,
    BPE.process(executors: e, docs: [doc | _]) = proc) do
    newDoc = :kvs.setfield(doc, :project, rejectComment(:kvs.field(doc, :project)))
    newProc = BPE.process(proc, docs: [newDoc])
    BPE.result(state: generalAction(r, proc, newProc), executed: e,
      continue: [
        continueGeneral(r, newProc, newDoc),
        BPE.continue(type: :stop)
      ])
  end
end
\end{verbatim}
\endgroup

\subsubsection{newOutputProc.ex (Процес обробки вихідних документів)}

\begingroup
\footnotesize
\begin{verbatim}
defmodule NewOutput.Proc do
  require FORM
  require BPE
  require ERP
  require CRM.KEP
  require KVS
  require General.Proc
  require Reports
  import General.Proc

  @sev_send_type '170'

  def def(), do: BPE.xml("/priv/bpmn/newOutputProc.bpmn")
  def auth(_roles), do: true

  def general_specific_statistic({:callback, _, _}, result, _), do: result
  def personal_specific_statistic({:callback, _, _}, result, _), do: result
  def general_statistic({:callback, _, _}, result, _), do: result
  def main_document_statistic({:callback, _, _}, result, _), do: result
  def personal_statistic({:callback, _, _}, result, _), do: result
  def remove_personal_statistic(_, _, _, _, _), do: []

  def history({:callback, _, t}, BPE.result(executed: executed,
    state: BPE.process(id: pid, parent: parent, docs: [doc | _],
    executors: executors)) = result, :created) do
    History.generateHistory(%HistExt{action: :created, procId: pid,
      user: executed, userTo: executors, stage: t, newValue: doc})
    History.generateHistory(%HistExt{action: :add_proc, procId: parent,
      subPid: pid, user: :kvs.field(doc, :registered_by)})
    result
  end

  def history({:callback, _, _}, BPE.result(state: BPE.process(id: pid,
    docs: [doc | _], parent: p)) = result, :familiarization) do
    History.generateHistory(%HistExt{action: :familiarization, procId: pid,
      newValue: doc})
    History.generateHistory(%HistExt{action: :familiarized, procId: p,
      newValue: doc, subPid: pid})
    result
  end

  def history({:callback, _, _}, BPE.result(executed: executed,
    state: BPE.process(id: pid, docs: [doc | _])) = result, action)
    when action in [:agreed, :signed, :certified, :execution] do
    History.generateHistory(%HistExt{action: action, procId: pid,
      user: executed, newValue: doc})
    result
  end

  def history({:callback, _, t}, BPE.result(executed: executed,
    state: BPE.process(id: pid, executors: executors)) = result, _) do
    History.generateHistory(%HistExt{action: :stage, procId: pid,
      user: executed, userTo: executors, stage: t})
    result
  end

  def report({:callback, _, _}, BPE.result(state: proc) = result, {action, t}) do
    :erlang.apply(Reports, action, [Reports.info(report_type: t, process: proc)])
    result
  end

  def action(BPE.broadcastEvent(name: 'ImplementationRejected'), proc) do
    BPE.result(state: proc, continue: [
      BPE.continue(fn: :add, type: :spawn, module: :'Elixir.Reports',
        args: [Reports.info(report_type: :expired, types: :implementation,
        process: proc)])
    ])
  end

  def action(BPE.broadcastEvent(name: 'Delete',
    sender: BPE.process(module: m)) = ev, BPE.process(id: pid) = proc)
    when m in [NewInput.Proc, NewOrg.Proc, NewInternal.Proc] do
    :bpe.broadcastEvent(pid, ev)
    BPE.result(state: proc, continue: [
      BPE.continue(type: :stop),
      BPE.continue(type: :spawn, fn: :delete_process,
        module: :"Elixir.General.Proc", args: [pid])
    ])
  end

  def action(BPE.broadcastEvent(name: 'Delete') = ev,
    BPE.process(id: pid) = proc) do
    :bpe.broadcastEvent(pid, ev)
    BPE.result(type: :stop, state: proc)
  end

  def action(BPE.asyncEvent(name: 'SetStatus', payload: {group, s, user}),
    BPE.process(id: pid, module: m, executors: execs, docs: [doc | _]) = proc) do
    request = case :bpe.head(pid) do
      BPE.hist(task: BPE.sequenceFlow(source: src, target: dst)) ->
        {:request, src, dst}
      _ -> []
    end
    newProc = BPE.process(proc, executors: setExecutorStatus(execs, user, s))
    m.generalMoveTo(request, newProc, doc, "#{group}")
    BPE.result(state: newProc, continue: [
      continueGeneral(request, newProc, doc, group),
      BPE.continue(fn: :generateHistory, type: :spawn, module: :'Elixir.History',
        args: [%HistExt{action: s, procId: pid, user: user}])
    ])
  end
end
\end{verbatim}
\endgroup

\subsubsection{abacPDP.ex (Модуль прийняття рішень доступу)}

\begingroup
\footnotesize
\begin{verbatim}
defmodule ABAC.PDP do
  require ABAC
  require KVS

  def get_policy([ABAC.policy() = h | t]), do: get_policy(h) ++ get_policy(t)
  def get_policy(ABAC.policy(object: o, api_endpoint: e)) do
    :lists.filter(fn ABAC.policy(api_endpoint: ^e) -> true; _ -> false end,
      :kvs.index(:policy, :object, o, KVS.kvs(mod: :kvs_mnesia)))
  end
  def get_policy(_), do: []

  def decision(req0) do
    {r, p} = ABAC.PIP.request(req0)
    :lists.any(fn x -> :lists.any(&policy(r, &1), get_policy(x)) end, p)
  end

  def deny(req0) do
    {r, p} = ABAC.PIP.request(req0)
    :lists.any(fn x -> :lists.any(&!policy(r, &1), get_policy(x)) end, p)
  end

  def get(req0) do
    {ABAC.request() = req, policy} = ABAC.PIP.request(req0)
    :lists.flatten(:lists.map(fn x ->
      :lists.filtermap(fn
        ABAC.rule(values: []) -> false
        ABAC.rule(values: v, object: o) -> {true, :kvs.field(o, v)}
      end, policy(req, x))
    end, get_policy(policy)))
  end

  def parse(r), do: ABAC.PIP.parse(r)

  def policy(ABAC.request(type: t) = req, ABAC.policy(rules: r, combining: c))
    when t in [:decision, :deny] do
    :erlang.apply(:lists, c, [
      fn ref -> rule(req, :erlang.element(2,
        :kvs.get(:rule, ABAC.rule_ref(ref, :id), KVS.kvs(mod: :kvs_mnesia)))) end,
      r
    ])
  end
  def policy(ABAC.request(type: :get) = req, ABAC.policy(rules: r,
    combining: :all)) do
    :erlang.element(2, :lists.foldr(fn
      _, {false, _} = x -> x
      ABAC.rule_ref(id: i), {true, acc} ->
        rule = :erlang.element(2, :kvs.get(:rule, i, KVS.kvs(mod: :kvs_mnesia)))
        if rule(req, rule), do: {true, [rule | acc]}, else: {false, []}
    end, {true, []}, r))
  end
  def policy(ABAC.request(type: :get) = req, ABAC.policy(rules: r)) do
    :lists.foldr(fn ABAC.rule_ref(id: i), acc ->
      x = :erlang.element(2, :kvs.get(:rule, i, KVS.kvs(mod: :kvs_mnesia)))
      case rule(req, x) do
        true -> [x | acc]
        _ -> acc
      end
    end, [], r)
  end
  def policy(ABAC.request(type: :decision), _), do: false
  def policy(ABAC.request(type: :deny), _), do: true
  def policy(ABAC.request(type: :get), _), do: []

  def rule(ABAC.request(endpoint: e, resources: resources) = req,
    ABAC.rule(api_endpoint: e, type: t, object_condition: [], object: o,
    subject_condition: c, resource_match: rm) = r) do
    x = :lists.filter(fn ^o -> true; _ -> false end, resources)
    result(t, x != [] and :erlang.apply(:lists, rm, [
      fn _ -> :erlang.apply(:"Elixir.ABAC.Condition", c, [req, r]) end,
      x
    ]))
  end
  def rule(ABAC.request(endpoint: e, resources: resources) = req,
    ABAC.rule(api_endpoint: e, type: t, object_condition: oc, object: o,
    subject_condition: c, resource_match: rm) = r) do
    x = :lists.filter(&(:erlang.is_record(&1, :erlang.element(1, o))), resources)
    result(t, x != [] and :erlang.apply(:lists, rm, [
      fn item -> :erlang.apply(:"Elixir.ABAC.Condition", c, [req, r]) and
        :erlang.apply(:"Elixir.ABAC.Condition", oc, [item, r]) end,
      x
    ]))
  end
  def rule(_, _), do: false

  defp result(:permit, x), do: x
  defp result(:deny, x), do: !x
end
\end{verbatim}
\endgroup

\subsubsection{kep\_validator.erl (Модуль валідації електронних підписів)}

\begingroup
\footnotesize
\begin{verbatim}
-module(kep_validator).
-export([validate_signature/2, verify_time_stamp/1]).

validate_signature(FileHash, SignatureBytes) ->
    case ca_client:verify_signature(FileHash, SignatureBytes) of
        {ok, SignerCert} ->
            case ca_client:check_cert_status(SignerCert) of
                valid -> ok;
                expired -> {error, cert_expired};
                revoked -> {error, cert_revoked}
            end;
        Error -> Error
    end.

verify_time_stamp(TimeStampToken) ->
    case ca_client:verify_tsp(TimeStampToken) of
        true -> ok;
        false -> {error, invalid_tsp}
    end.
\end{verbatim}
\endgroup

\subsection{Інсталяція та конфігурація системи (mix.exs)}
Запуск та збирання релізу системи електронного документообігу на базі Erlang/Elixir здійснюється наступним чином:

\subsubsection{Пререквізити для операційної системи (Ubuntu)}
\begin{verbatim}
$ sudo apt install erlang elixir build-essential libcsv3 libcsv-dev cmake unoconv
$ sudo snap install onlyoffice-ds
\end{verbatim}

\subsubsection{Залежності проекту (mix.exs)}
\begin{verbatim}
defmodule CRM.Mixfile do
  use Mix.Project
  def project do
    [ app: :crm, version: "3.1.0", elixir: "~> 1.9", deps: deps() ]
  end
  def deps do
    [
      {:cowboy, "~> 2.9.0"},
      {:cowlib, "~> 2.11.0"},
      {:rocksdb, "~> 1.6.0"},
      {:sax, "~> 1.0.0"},
      {:ca, "~> 1.11.0"},
      {:barlix2, "~> 0.6.0"},
      {:qrcode, "~> 0.1"},
      {:n2o, "~> 8.12.1", override: true},
      {:soa, "~> 0.1.7"},
      {:jsone, "~> 1.5.1"},
      {:bpe, "7.6.3"},
      {:kvs, "~> 9.7.0", override: true},
      {:smtp, "~> 1.1.2"},
      {:pop3, "~> 2.1.0"},
      {:nitro, "~> 7.7.1", override: true},
      {:form, "~> 7.7.0"}
    ]
  end
end
\end{verbatim}

\subsubsection{Компіляція та запуск у режимі демона}
\begin{verbatim}
$ mix deps.get
$ mix release
$ _build/dev/rel/crm/bin/crm daemon
\end{verbatim}

\newpage
\section{Вимоги до засобу за ISO/IEC/IEEE 29148}

\subsection{Функціональні вимоги}
\begin{itemize}
    \item \textbf{[REQ-MAIL-FUN-001]} Модуль «Документи» повинен автоматично нумерувати вхідні, вихідні та внутрішні документи відповідно до заданих шаблонів нумерації та календарного року.
    \item \textbf{[REQ-MAIL-FUN-002]} Модуль «Документи» повинен наносити QR-код (21x21 мм) та лінійний штрих-код на першу сторінку документа під час реєстрації.
    \item \textbf{[REQ-MAIL-FUN-003]} Модуль «Документи» повинен виконувати перевірку кваліфікованих електронних підписів (КЕП) та печаток (КЕПч) за протоколами OCSP/CRL.
    \item \textbf{[REQ-MAIL-FUN-004]} Модуль «Документи» повинен підтримувати паралельне виконання завдань (підзадач) співвиконавцями в межах однієї РМК.
    \item \textbf{[REQ-MAIL-FUN-005]} Модуль «Документи» повинен здійснювати автоматичний контроль термінів виконання та надсилати сповіщення при наближенні дати прострочення.
    \item \textbf{[REQ-MAIL-FUN-006]} Модуль «Документи» повинен забезпечувати інтеграцію з СЕВ ОВВ для відправки та прийому документів через конверти MHS X.420.
    \item \textbf{[REQ-MAIL-FUN-007]} Модуль «Документи» повинен дозволяти делегування прав керівника помічнику або заміснику із фіксацією періоду делегування.
\end{itemize}

\subsection{Вимоги до інтерфейсів та дизайну}
\begin{itemize}
    \item \textbf{[REQ-MAIL-UI-001]} Користувацький інтерфейс повинен відповідати вимогам WCAG 2.1 рівню AA щодо вебдоступності.
    \item \textbf{[REQ-MAIL-UI-002]} Інтерфейс має бути побудований за технологією Single Page Application (SPA) на базі WebSocket з'єднань.
\end{itemize}

\subsection{Вимоги до безпеки та захисту інформації}
\begin{itemize}
    \item \textbf{[REQ-MAIL-SEC-001]} Авторизація користувачів повинна відбуватися виключно за допомогою сертифікатів КЕП X.509.
    \item \textbf{[REQ-MAIL-SEC-002]} Доступ до будь-які РМК повинен перевірятися підсистемою ABAC на основі атрибутів користувача та об'єкта.
    \item \textbf{[REQ-MAIL-SEC-003]} Вся взаємодія між клієнтом та сервером повинна здійснюватися через шифроване з'єднання TLS 1.3.
\end{itemize}

\subsection{Вимоги до продуктивності}
\begin{itemize}
    \item \textbf{[REQ-MAIL-PERF-001]} Час відкриття та відображення сторінки РМК не повинен перевищувати 500 мс.
    \item \textbf{[REQ-MAIL-PERF-002]} Система повинна підтримувати обробку не менше 500 паралельних WebSocket з'єднань на один процесорний вузол.
\end{itemize}

\subsection{Вимоги до логування, тестування та супроводу}
\begin{itemize}
    \item \textbf{[REQ-MAIL-LOG-001]} Усі дії користувачів над РМК (створення, редагування, підписання) повинні логуватися в системному Audit Log у форматі JSON.
    \item \textbf{[REQ-MAIL-ADM-001]} Покриття вихідного коду автоматичними тестами має бути не менше 80\%.
\end{itemize}

\end{document}
